\clearpage\phantomsection
\addcontentsline{toc}{chapter}{Resumen y Palabras Clave}
\vspace*{0.8cm}

% ---------- RESUMEN (SPANISH) — max 500 words ----------
\Large{\textsc{\textcolor{nar}{Resumen}}} \par
\vspace{0.2cm}
\normalsize{\selectlanguage{spanish}
La enfermedad de Parkinson es el segundo trastorno neurodegenerativo más
frecuente, y la dificultad para caminar figura entre sus signos más tempranos
e incapacitantes. El radar de onda continua modulada en frecuencia (FMCW)
permite observar la marcha sin contacto, sin cámaras y sin sensores
corporales: la grabación es una \textbf{firma micro-Doppler}, una imagen
tiempo-frecuencia que integra el movimiento de todas las partes del cuerpo y
de la que pueden extraerse las velocidades del torso y de las extremidades.
Por ello el radar resulta atractivo para evaluar el movimiento en la clínica
y en el hogar; sin embargo, hasta donde sabemos, ningún trabajo publicado
había evaluado un clasificador de Parkinson frente a controles sobre estos
firmas con \textbf{validación independiente del sujeto}, aquella en
la que cada persona puntuada es alguien que el modelo nunca vio durante el
entrenamiento, que es la condición real de una herramienta de cribado.

Este trabajo construye esa evaluación sobre una red de radar clínicamente
validada y una cohorte de pacientes de Parkinson y controles sanos. Se
desarrolla un proceso completo de análisis y se evalúan desde modelos
clásicos hasta varias redes neuronales profundas, incluida la arquitectura
de referencia de la literatura, con validación cruzada dejando un sujeto
fuera; cada modelo se compara con referencias sencillas, como un predictor
que solo conoce la duración de cada grabación, y cada resultado se acompaña
de su intervalo de confianza.

Bajo esta evaluación, ningún clasificador demuestra conocer la enfermedad
más allá de la duración de la grabación: los pacientes sencillamente tardan
más. El análisis de atención atribuye la puntuación restante a propiedades
de cada grabación y no a la marcha; y reproducir sobre los mismos datos el
reparto aleatorio habitual en la literatura muestra cuánto infla ese diseño
el rendimiento: modelos aparentemente listos para la clínica se desploman
ante desconocidos. Lo que sobrevive es clínicamente coherente: el radar mide
con fiabilidad los tiempos de la marcha, una prueba \textit{Timed-Up-and-Go}
automatizada.

La contribución es doble: la primera evaluación independiente del sujeto de
este problema, y una demostración de cómo los clasificadores de marcha por
radar llegan a parecer mejores de lo que son, junto con el protocolo de
evaluación que lo impide.
\selectlanguage{english}}
\vspace{0.9cm}

% ---------- SUMMARY (ENGLISH) — max 500 words ----------
\Large{\textsc{\textcolor{nar}{Summary}}} \par
\vspace{0.2cm}
\normalsize{%
Parkinson's disease is the second most common neurodegenerative disorder,
and difficulty walking is among its earliest and most disabling signs.
Frequency-Modulated Continuous Wave (FMCW) radar can observe a walk without
contact, cameras or wearable sensors: the recording is a
\textbf{micro-Doppler signature}, a time-frequency image that integrates the
motion of every part of the body, from which the velocities of the torso and
the limbs can be extracted. Radar is therefore attractive for assessing
movement in clinics and homes, yet, to the best of our knowledge, no
published work had evaluated a Parkinson's-versus-control classifier on
these signatures under \textbf{subject-independent validation}, in which
every scored person is one the model never saw during training, which is the
condition a screening tool faces in practice.

This thesis builds that evaluation on a clinically validated radar network
and a cohort of Parkinson's patients and healthy controls. A full pipeline
is developed, and models from classical baselines to several deep neural
networks, including the field's reference architecture, are evaluated with
leave-one-subject-out cross-validation; every model is compared against
simple baselines, such as a predictor that knows only how long each
recording lasts, and every result carries a confidence interval.

Under this evaluation, no classifier demonstrates more knowledge of the
disease than recording duration alone: patients simply take longer.
Attention analysis traces the remaining score to properties of each
recording rather than of the walk, and reproducing the literature's usual
random hold-out on the same data shows how strongly that design inflates
performance: models that look clinic-ready on familiar subjects collapse on
strangers. What survives is clinically coherent: radar reliably measures
the timing of the walk, an automated \textit{Timed-Up-and-Go} test.

The contribution is twofold: the first subject-independent evaluation of
this problem, and a demonstration of how radar gait classifiers come to
look better than they are, together with the evaluation protocol that
prevents it.}
\vspace{0.9cm}

% ---------- KEYWORDS ----------
\Large{\textsc{\textcolor{nar}{Palabras Clave}}} \par
\vspace{0.2cm}
\normalsize{Enfermedad de Parkinson; radar FMCW; micro-Doppler; análisis de la marcha;
aprendizaje profundo; validación independiente del sujeto; control de factores de confusión.} \par
\vspace{0.8cm}

\Large{\textsc{\textcolor{nar}{Keywords}}} \par
\vspace{0.2cm}
\normalsize{Parkinson's disease; FMCW radar; micro-Doppler; gait analysis;
deep learning; subject-independent validation; confound control.} \par
